Basta meter una cuchara en un vaso con agua para notar que algo raro pasa con la luz. La cuchara parece quebrada justo donde entra al líquido.

Esto tiene nombre: refracción.

Cuando la luz pasa de un material a otro, por ejemplo del aire al vidrio, cambia su velocidad, y ese cambio hace que el rayo se doble en la frontera entre los dos medios; qué tanto se dobla depende de cada material y del ángulo con el que llega la luz, y la regla que lo describe se conoce como ley de Snell, en honor al matemático holandés Willebrord Snell, que la escribió a principios del siglo XVII aunque otros ya la habían intuido antes.

En la práctica anterior vimos que esta ley se puede deducir del principio de Fermat. Ahora tocaba verla en acción. Con una fuente de luz, una mesa giratoria graduada en grados y una lente de acrílico con forma de media luna, medimos con qué ángulo salía el rayo para distintos ángulos de entrada, tanto cuando la luz pasaba del aire al acrílico como cuando hacía el viaje contrario, y con esos datos calculamos el índice de refracción del acrílico y su ángulo crítico.
